\section{During Measurements}
After all preparations have been completed, a measurement procedure has been designed, and the required instruments have been selected, the finalized measurement can begin. Ideally, the preparations are so well defined that one simply follows the designed procedure.

However, this is not always the case. This section describes what should be done during measurements and highlights key considerations to be aware of.

\subsection{Records}
Once the measurement begins, the time and place should be recorded, along with the name of the operator, the instruments used, and ideally their serial numbers. Metrology-grade measurements also require recording environmental conditions, such as temperature and humidity. While this may not be strictly necessary for typical experimental work, it is important to be aware of their influence: most laboratory instruments are designed to operate around 23$^\circ$C and 45\% relative humidity.

The environment in which the measurements are performed should always be described. All measurement data should be recorded in accordance with the established measurement procedure.

\subsection{Stability}
During measurements, the experimenter should monitor the system and data for stability, drift, or signs that an instrument may be saturated or operating out of range. If such issues occur, they must be noted, and an evaluation should be made to determine whether the measurements are still acceptable or if the procedure, measurement ranges, or instruments need to be re-assessed. These observations and considerations should be documented and discussed with the supervisor.